% Generated draft: requires mathematical and source review.
\documentclass[11pt]{article}
\usepackage[margin=25mm]{geometry}
\usepackage{fontspec}
\setmainfont{Latin Modern Roman}
\usepackage{amsmath,amssymb,mathtools,unicode-math}
\setmathfont{Latin Modern Math}
\usepackage{enumitem,needspace,xurl,hyperref}
\setlength{\parindent}{0pt}
\setlength{\parskip}{5pt}
\setlength{\emergencystretch}{3em}
\newcommand{\N}{\mathbb{N}}
\newcommand{\Z}{\mathbb{Z}}
\newcommand{\Q}{\mathbb{Q}}
\newcommand{\R}{\mathbb{R}}
\newcommand{\C}{\mathbb{C}}
\newcommand{\F}{\mathbb{F}}
\newcommand{\A}{\mathbb{A}}
\newcommand{\PP}{\mathbb{P}}
\newcommand{\E}{\mathbb{E}}
\newcommand{\eps}{\varepsilon}
\newcommand{\dd}{\,\mathrm{d}}
\begin{document}
\section*{UnsolvedMath ID: 30005621}
\textbf{Source problem code:} catalogue entry
\subsection{Definitions and mathematical statement}
The supplied record introduces
\[
\mathcal R_n=\{P:P\text{ is a polygon with exactly }n\text{ sides}\},
\]
where \(n\) is fixed. The original extract does not specify the ambient space, admissibility conditions on polygons, or the range of \(n\). Under the usual interpretation as nondegenerate planar polygons, one would require an integer \(n\geq 3\); this interpretation is not an explicit hypothesis in the extract.

The reconstructed statement supplied by the registry is: for polygons with a fixed number \(n\) of sides, optimizers of normalized Dirichlet and Neumann eigenvalues converge to the regular \(n\)-gon as the eigenvalue index tends to infinity.

This is an incomplete formulation rather than a fully specified mathematical conjecture. It does not identify the differential operator, the normalization of the eigenvalues, whether each spectral problem is minimized or maximized, or the admissible polygon class beyond the number of sides. It also does not specify a topology of convergence, how translation, rotation and scale are handled, whether optimizers exist at every index, or whether convergence is required for every choice of optimizers or only for some selection. Consequently, no definite quantified convergence assertion can be extracted without additional hypotheses.

The original extract stops after introducing \(\mathcal R_n\) and then contains an apparently separate fragment headed ``Richard Laugesen: Neumann and Robin eigenvalues.'' That fragment mentions a domain \(\Omega\), a Robin parameter \(\alpha\), and a corrupted Rayleigh quotient. It supplies no complete polygon optimization claim and does not establish that Robin eigenvalues belong to the intended convergence problem.
\subsection{Short English statement}
The supplied reconstructed claim concerns convergence of high-index Dirichlet and Neumann spectral optimizers among polygons with a fixed number of sides to a regular polygon. The normalization, optimization directions, admissible domains and meaning of convergence are missing, so a precise conjecture cannot yet be stated. The registry labels the problem open; this is a supplied status, not an independently verified conclusion.
\subsection{Sources}
\begin{itemize}
\item Supplied local registry record, permanent UnsolvedMath ID \(30005621\), title ``High-Index Spectral Optimizers Among Polygons,'' exact source problem code \texttt{catalogue entry}, in the Oberwolfach Reports Open Problems collection.
\item Supplied source citation: \emph{Geometric Spectral Theory (2024)}, Oberwolfach Reports, DOI \texttt{10.4171/owr/2023/36}, pp.\ 2017--2106. Source URL: \url{https://doi.org/10.4171/owr/2023/36}. The record gives proposed year \(2023\), with no named proposer.
\item Supplied literature source, labelled ``original OWR report'': \url{https://doi.org/10.4171/owr/2023/36}.
\item Literature checked at: \texttt{2026-08-22}. This is supplied dated curation, not independently verified. The supplied assessment labels the problem open and says that the extracted record stops after defining the class of \(n\)-sided polygons, leaving no specific extremal claim whose resolution can be identified.
\item The registry separately labels its reconstructed statement \texttt{corrected\_verified} and reports that it was checked against the official Oberwolfach report. These are supplied verification claims; no independent source consultation is asserted here. The adjacent Richard Laugesen heading is retained as fragment attribution, not as attribution of the polygon problem.
\end{itemize}
\end{document}
